% Part: turing-machines
% Chapter: machines-computations
% Section: turing-machines

\documentclass[../../../include/open-logic-section]{subfiles}

\begin{document}

\olfileid{tur}{mac}{tur}
\olsection{ماشین‌های تورینگ}

\begin{explain}
تعریف رسمیِ آنچه ماشین تورینگ را می‌سازد انتزاعی به نظر می‌رسد، اما
در واقع ساده است: فقط همهٔ اطلاعات لازم برای مشخص‌کردن کار ماشین
تورینگ را در یک ساختار ریاضی می‌گنجاند. این اطلاعات عبارت‌اند از:
(۱)~حالت‌هایی که ماشین می‌تواند در آن‌ها باشد، (۲)~نمادهای مجاز روی
نوار، (۳)~حالتی که ماشین باید در آن آغاز کند، و (۴)~مجموعهٔ دستورهای
ماشین.
\end{explain}

\begin{defn}[ماشین تورینگ]
\emph{ماشین تورینگ} $M$ چهارتاییِ
$\langle Q,\Sigma,q_0,\delta\rangle$ است که از این اجزا تشکیل می‌شود:
\begin{enumerate}
\item مجموعهٔ متناهیِ \emph{حالت‌ها}~$Q$؛
\item \emph{الفبای} متناهیِ $\Sigma$ که $\TMendtape$ و~$\TMblank$ را
  دربر می‌گیرد؛
\item \emph{حالت آغازین}~$q_0\in Q$؛
\item \emph{مجموعهٔ دستورهای} متناهی
  $\delta\colon Q\times\Sigma
  \pto Q\times\Sigma\times\{\TMleft,\TMright,\TMstay\}$.
\end{enumerate}
تابع جزئیِ~$\delta$ را \emph{تابع گذارِ}~$M$ نیز می‌نامند.
\end{defn}

\begin{explain}
فرض می‌کنیم نوار تنها در یک جهت نامتناهی است. به همین دلیل مفید است
نماد خاص~$\TMendtape$ را نشانگر انتهای چپ نوار قرار دهیم. این کار
تشخیص زمانی را که برنامهٔ ماشین تورینگ «در خطر» بیرون‌رفتن از نوار
است آسان‌تر می‌کند. می‌توانستیم فرض کنیم این نماد هرگز بازنویسی
نمی‌شود؛ یعنی اگر $\delta(q,\TMendtape)$ تعریف شده باشد، آنگاه
$\delta(q,\TMendtape)=\tuple{q',\TMendtape,D}$. بعضی کتاب‌های درسی
چنین می‌کنند، اما ما نمی‌کنیم. هنگام ساخت ماشین تورینگتان صرفاً
مراقب باشید هرگز $\TMendtape$ را بازنویسی نکند. افزون بر این، گاه
اجازه‌دادن به چنین بازنویسی‌ای انعطاف سودمندی فراهم می‌کند.
\end{explain}

\begin{ex}
\emph{ماشین زوج:} ماشین زوج به‌طور رسمی چهارتایی
$\tuple{Q,\Sigma,q_0,\delta}$ است که در آن
\begin{align*}
Q & = \{ q_0, q_1 \} \\
\Sigma & = \{ \TMendtape, \TMblank, \TMstroke \}, \\
\delta(q_0, \TMstroke) & = \tuple{q_1, \TMstroke, \TMright},\\
\delta(q_1, \TMstroke) & = \tuple{q_0, \TMstroke, \TMright},\\
\delta(q_1, \TMblank)  & = \tuple{q_1, \TMblank, \TMright}.
\end{align*}
\end{ex}

\end{document}
